\chapter{대결 관리}

\begin{summarybox}{한눈에}
학생이 다른 학생·AI 와 \textbf{1:1 듀얼 대결}로 학습을 즐기는 시스템. 본사가 시즌과 AI 플레이어를 운영하며, 기관 관리자는 본 기관 학생의 \textbf{매치 / 시즌 점수 / 보상 결과}를 모니터링합니다.
\end{summarybox}

\kfShot{10-duel}{대결 관리 — 본 기관 학생 듀얼 현황}

\section*{화면 구조}

\screentabs{\activetab{매치 현황} \quad \inacttab{시즌 점수} \quad \inacttab{보상 / 신고}}

\begin{screenbox}{본 기관 학생 듀얼 활동}
\noindent
\begin{tabular}{@{}p{2cm}p{2.5cm}p{2cm}p{2cm}p{2cm}p{2cm}@{}}
\toprule
\textbf{시즌} & \textbf{학생} & \textbf{매치 수} & \textbf{승} & \textbf{패} & \textbf{현재 연승} \\
\midrule
2026 봄 & 홍길동 & 24 & 14 & 10 & 3 \\
2026 봄 & 김영희 & 18 & 12 & 6 & 5 \\
\bottomrule
\end{tabular}
\end{screenbox}

\section*{여기서 할 수 있는 일}
\begin{itemize}[leftmargin=1.2em]
  \item 본 기관 학생들의 \textbf{매치 이력} 조회 (상대·점수·결과)
  \item \textbf{시즌 점수 / 랭킹} 확인
  \item 학생 듀얼 활동을 통합 성적표와 연결해서 보기
  \item 부정행위·욕설 \textbf{신고 처리} (본사 단독 신고는 본사가 처리)
\end{itemize}

\section*{자주 쓰는 흐름}
\begin{enumerate}
  \item ``매치 현황'' → 본 기관 학생 듀얼 활동 점검
  \item 학생 클릭 → 상세 매치 로그 (어떤 문제·어떤 시간에·결과)
  \item 학습 계획표와 함께 보면서 ``듀얼이 학습으로 이어지고 있는가'' 판단
\end{enumerate}

\begin{tipbox}{알아두면 좋은 점}
\begin{itemize}[leftmargin=1.2em]
  \item 시즌은 본사가 발표·종료 — 본 기관이 별도 시즌을 만들 수는 없음
  \item AI 플레이어 슬라이더(난이도 자동 조정) 는 본사가 설정 — 학생별 매칭 난이도가 자동으로 맞춰짐
  \item 신고 처리 후 학생 일시 정지·시즌 점수 차감은 본사 관리자가 최종 결정
\end{itemize}
\end{tipbox}
